\documentclass[UTF8,a4paper,12pt]{ctexart}

\usepackage{amsmath,amssymb}
\usepackage{geometry}
\usepackage{booktabs}
\usepackage{array}
\usepackage{siunitx}

\geometry{left=2.5cm,right=2.5cm,top=2.5cm,bottom=2.5cm}
\sisetup{per-mode=symbol}

\title{第一问建模思路与公式整理}
\author{}
\date{}

\begin{document}
\maketitle

\section{建模简化与总体思路}

根据手写稿，Hussain--Dincer（2003）给出了轴对称二维圆柱模型：
\begin{equation}
    T=T(r,z,t),\qquad m=m(r,z,t),
\end{equation}
其中热传导和水分扩散同时包含径向与轴向过程，并在圆柱侧壁以及上下端设置对流边界条件。

手写稿进一步引用 de Silva（2014）的处理：当圆柱长度远大于半径，即
\begin{equation}
    L\gg R,
\end{equation}
时，可将二维轴对称问题近似化为一维径向圆柱扩散问题。

题中棒材的几何尺度记为
\begin{equation}
    L=\SI{25}{cm},\qquad R=\SI{2}{cm}.
\end{equation}
由于研究重点为棒材不同径向位置处的温度和含水率变化，因此忽略端部效应，只考虑径向传热、传质过程，令
\begin{equation}
    T=T(r,t),\qquad C=C(r,t),\qquad 0\le r\le R.
\end{equation}

第一问的物理过程可概括为：
\[
\text{热空气}\longrightarrow\text{棒材表面}\longrightarrow\text{棒材内部},
\]
其传热过程满足 Fourier 定律；水分迁移过程可概括为
\[
\text{棒材内部水分}\longrightarrow\text{棒材表面}\longrightarrow\text{空气},
\]
其传质过程满足 Fick 定律。

\section{温度场模型}

\subsection{控制方程}

取半径为 $r$、厚度为 $\mathrm{d}r$ 的圆柱环状微元。径向热流密度由 Fourier 定律给出：
\begin{equation}
    q_r=-k\frac{\partial T}{\partial r},
\end{equation}
其中 $k$ 为材料导热系数。

由圆柱坐标系下的能量守恒，可得一维径向非稳态导热方程
\begin{equation}
    \rho C_p\frac{\partial T}{\partial t}
    =
    \frac{1}{r}\frac{\partial}{\partial r}
    \left(
        r k\frac{\partial T}{\partial r}
    \right),
    \qquad 0<r<R.
    \label{eq:heat-pde}
\end{equation}

\subsection{初始条件与边界条件}

初始温度为
\begin{equation}
    T(r,0)=28^\circ\mathrm{C}.
\end{equation}

在圆柱中心 $r=0$ 处，由轴对称性有
\begin{equation}
    \left.\frac{\partial T}{\partial r}\right|_{r=0}=0.
    \label{eq:heat-center-bc}
\end{equation}

在棒材外表面 $r=R$ 处采用对流换热边界条件：
\begin{equation}
    -k\left.\frac{\partial T}{\partial r}\right|_{r=R}
    =
    h\left[T(R,t)-T_a(t)\right],
    \label{eq:heat-surface-bc}
\end{equation}
其中 $h$ 为题目给定的对流换热系数，$T_a(t)$ 为烘房空气温度。

手写稿记录了
\begin{equation}
    T_a(0)=28^\circ\mathrm{C},\qquad
    T_a(600)=33.202^\circ\mathrm{C}.
\end{equation}
对于其余时刻，可根据题目所给离散数据对 $T_a(t)$ 作分段线性插值。若
$t_j\le t\le t_{j+1}$，则
\begin{equation}
    T_a(t)
    =
    T_{a,j}
    +
    \frac{T_{a,j+1}-T_{a,j}}{t_{j+1}-t_j}
    (t-t_j).
    \label{eq:Ta-interp}
\end{equation}

\section{水分扩散模型}

\subsection{控制方程}

令 $C(r,t)$ 表示材料的干基含水率，单位为 $\mathrm{kg/kg}$。
根据 Fick 定律，径向水分通量为
\begin{equation}
    J_C=-D(C)\frac{\partial C}{\partial r},
\end{equation}
其中 $D(C)$ 为随含水率变化的有效扩散系数。

圆柱坐标系下的一维径向水分扩散方程为
\begin{equation}
    \frac{\partial C}{\partial t}
    =
    \frac{1}{r}\frac{\partial}{\partial r}
    \left[
        rD(C)\frac{\partial C}{\partial r}
    \right],
    \qquad 0<r<R.
    \label{eq:moisture-pde}
\end{equation}

手写稿给出的扩散系数模型为
\begin{equation}
    D(C)=7\times 10^{-9}\exp(-0.89C).
    \label{eq:Dc}
\end{equation}

\subsection{初始条件与边界条件}

初始含水率为
\begin{equation}
    C(r,0)=2.55.
\end{equation}

圆柱中心满足对称边界条件
\begin{equation}
    \left.\frac{\partial C}{\partial r}\right|_{r=0}=0.
    \label{eq:moisture-center-bc}
\end{equation}

外表面采用对流传质边界条件
\begin{equation}
    -D(C)
    \left.\frac{\partial C}{\partial r}\right|_{r=R}
    =
    h_m\left[C(R,t)-C_a(t)\right],
    \label{eq:moisture-surface-bc}
\end{equation}
其中手写稿给出
\begin{equation}
    h_m=8\times 10^{-7}\ \mathrm{m/s},
\end{equation}
$C_a(t)$ 表示烘房空气中的水分浓度（由附件~1 给出的数据确定）。若附件中给出的是离散时刻数据，同样可采用分段线性插值得到连续函数 $C_a(t)$。

\section{数值求解：有限体积法}

\subsection{空间离散}

沿径向将区间 $[0,R]$ 划分为若干控制体。对第 $i$ 个控制体，
\begin{equation}
    r_{i-\frac12}\le r\le r_{i+\frac12}.
\end{equation}
其体积为
\begin{equation}
    V_i
    =
    \pi L
    \left(
        r_{i+\frac12}^2-r_{i-\frac12}^2
    \right),
    \label{eq:volume}
\end{equation}
左右两侧控制面的面积分别为
\begin{equation}
    A_{i-\frac12}=2\pi Lr_{i-\frac12},
    \qquad
    A_{i+\frac12}=2\pi Lr_{i+\frac12}.
    \label{eq:area}
\end{equation}

有限体积离散的核心思想是
\[
\boxed{\text{控制体内物理量变化率}
=
\text{左侧流入}
-
\text{右侧流出}}
\]
或等价地写成各控制面净通量之和。

\section{温度方程的有限体积离散}

\subsection{内部控制体}

对第 $i$ 个内部控制体进行能量守恒，有
\begin{equation}
    \rho C_pV_i\frac{\mathrm{d}T_i}{\mathrm{d}t}
    =
    kA_{i-\frac12}\frac{T_{i-1}-T_i}{\Delta r}
    +
    kA_{i+\frac12}\frac{T_{i+1}-T_i}{\Delta r}.
    \label{eq:heat-fvm-semidiscrete}
\end{equation}

令
\begin{equation}
    t_n=n\Delta t,
    \qquad
    T_i^n\approx T(r_i,t_n).
\end{equation}
时间方向采用全隐式格式，
\begin{equation}
    \frac{\mathrm{d}T_i}{\mathrm{d}t}
    \approx
    \frac{T_i^{n+1}-T_i^n}{\Delta t},
\end{equation}
并将右端空间通量全部取在第 $n+1$ 时间层，则
\begin{align}
    \rho C_pV_i
    \frac{T_i^{n+1}-T_i^n}{\Delta t}
    &=
    kA_{i-\frac12}
    \frac{T_{i-1}^{n+1}-T_i^{n+1}}{\Delta r}
    \notag\\
    &\quad+
    kA_{i+\frac12}
    \frac{T_{i+1}^{n+1}-T_i^{n+1}}{\Delta r}.
    \label{eq:heat-implicit}
\end{align}

整理为三对角形式：
\begin{equation}
    -a_iT_{i-1}^{n+1}
    +b_iT_i^{n+1}
    -c_iT_{i+1}^{n+1}
    =d_i,
    \label{eq:tridiag-row}
\end{equation}
其中
\begin{align}
    a_i&=\frac{kA_{i-\frac12}}{\Delta r},\\
    b_i&=\frac{\rho C_pV_i}{\Delta t}
        +\frac{k(A_{i-\frac12}+A_{i+\frac12})}{\Delta r},\\
    c_i&=\frac{kA_{i+\frac12}}{\Delta r},\\
    d_i&=\frac{\rho C_pV_i}{\Delta t}T_i^n.
\end{align}

因此每个时间步需要求解一个三对角线性方程组，其结构为
\begin{equation}
\begin{bmatrix}
b_1 & -c_1 & 0 & \cdots & 0\\
-a_2 & b_2 & -c_2 & \ddots & \vdots\\
0 & \ddots & \ddots & \ddots & 0\\
\vdots & \ddots & -a_{N-2} & b_{N-2} & -c_{N-2}\\
0 & \cdots & 0 & -a_{N-1} & b_{N-1}
\end{bmatrix}
\begin{bmatrix}
T_1^{n+1}\\
T_2^{n+1}\\
\vdots\\
T_{N-2}^{n+1}\\
T_{N-1}^{n+1}
\end{bmatrix}
=
\begin{bmatrix}
d_1\\
d_2\\
\vdots\\
d_{N-2}\\
d_{N-1}
\end{bmatrix}.
\end{equation}

\subsection{中心点 $r=0$ 的处理}

圆柱坐标方程在 $r=0$ 处形式上含有 $1/r$，因此直接代入会出现奇点。
采用中心控制体
\[
0\le r\le \frac{\Delta r}{2},
\]
其左侧控制面面积为
\begin{equation}
    A_{-\frac12}=0.
\end{equation}
于是
\begin{equation}
    \rho C_pV_0\frac{\mathrm{d}T_0}{\mathrm{d}t}
    =
    kA_{\frac12}
    \frac{T_1-T_0}{\Delta r}.
\end{equation}
其中
\begin{equation}
    V_0=\pi L\left(\frac{\Delta r}{2}\right)^2,
    \qquad
    A_{\frac12}=2\pi L\frac{\Delta r}{2}
    =\pi L\Delta r.
\end{equation}
从而得到
\begin{equation}
    \frac{\mathrm{d}T_0}{\mathrm{d}t}
    =
    \frac{4k}{\rho C_p(\Delta r)^2}
    (T_1-T_0),
    \label{eq:center-temp}
\end{equation}
这样即可避免直接处理 $1/r$ 奇点。

\subsection{外表面 $r=R$ 的处理}

令最外层节点为 $r_N=R$。外表面对流热流可写为
\begin{equation}
    \dot q
    =
    h\left[T_a(t)-T_N\right].
\end{equation}
因此最外层控制体满足
\begin{equation}
    \rho C_pV_N\frac{\mathrm{d}T_N}{\mathrm{d}t}
    =
    kA_{N-\frac12}
    \frac{T_{N-1}-T_N}{\Delta r}
    +
    hA_R\left[T_a(t)-T_N\right],
    \label{eq:outer-temp}
\end{equation}
其中
\begin{equation}
    A_R=2\pi RL.
\end{equation}

\section{水分方程的有限体积离散}

对式~\eqref{eq:moisture-pde}采用同样的有限体积思想。第 $i$ 个内部控制体满足
\begin{equation}
    V_i\frac{\mathrm{d}C_i}{\mathrm{d}t}
    =
    D_{i-\frac12}A_{i-\frac12}
    \frac{C_{i-1}-C_i}{\Delta r}
    +
    D_{i+\frac12}A_{i+\frac12}
    \frac{C_{i+1}-C_i}{\Delta r}.
    \label{eq:moisture-fvm}
\end{equation}

界面扩散系数按手写稿采用相邻节点扩散系数的算术平均：
\begin{equation}
    D_{i+\frac12}
    =
    \frac{D(C_i)+D(C_{i+1})}{2}.
    \label{eq:D-interface}
\end{equation}

在最外层 $r=R$ 处，由
\begin{equation}
    -D(C)\frac{\partial C}{\partial r}
    =
    h_m(C_s-C_a),
\end{equation}
得到外层控制体的半离散方程
\begin{equation}
    V_N\frac{\mathrm{d}C_N}{\mathrm{d}t}
    =
    D_{N-\frac12}A_{N-\frac12}
    \frac{C_{N-1}-C_N}{\Delta r}
    -
    h_mA_R(C_N-C_a).
    \label{eq:outer-moisture}
\end{equation}

由于 $D=D(C)$ 随含水率变化，水分方程属于非线性扩散方程。实际计算时可在每一时间步根据当前（或迭代更新后的）含水率重新计算 $D(C_i)$ 与 $D_{i+\frac12}$。

\section{第一问的计算流程}

根据上述模型，第一问可按以下流程求解：
\begin{enumerate}
    \item 根据附件数据构造连续的烘房空气温度 $T_a(t)$ 与空气水分浓度 $C_a(t)$；
    \item 将棒材半径区间 $[0,R]$ 划分为有限个径向控制体；
    \item 利用全隐式有限体积格式逐时间步求解温度场 $T(r,t)$；
    \item 根据 $D(C)=7\times10^{-9}e^{-0.89C}$ 更新有效水分扩散系数，并求解含水率场 $C(r,t)$；
    \item 输出题目要求的不同径向位置、不同时间处的温度与含水率，并据此完成第一问分析。
\end{enumerate}

% =========================
% 手写识别说明：
% 1. 文献在图片中仅出现“Hussain-Dincer 2003”和“de Silva 2014”，
%    未给出完整题名、期刊等信息，因此本文未补造参考文献条目。
% 2. 图片中的关键数值识别为：
%    L=25 cm, R=2 cm, T(r,0)=28°C, T_a(0)=28°C,
%    T_a(600)=33.202°C, C(r,0)=2.55,
%    D(C)=7×10^{-9} exp(-0.89 C), h_m=8×10^{-7} m/s。
% 3. “C(r,t)”在手写稿中表述为含水率（kg/kg）；本文按“干基含水率”整理。
% 4. 水分方程手写稿给出了空间有限体积半离散式与外边界处理；
%    温度方程明确写出了全隐式时间离散和三对角结构。
% =========================

\end{document}
